\documentclass[tikz,border=0pt]{standalone}

\usepackage{fix-cm}
\usepackage{fontspec}
\usetikzlibrary{arrows.meta,calc}

\setmainfont{Arial}
\setsansfont{Arial}

\definecolor{canvas}{HTML}{F7F9FC}
\definecolor{card}{HTML}{FFFFFF}
\definecolor{ink}{HTML}{172033}
\definecolor{muted}{HTML}{53627A}
\definecolor{border}{HTML}{C9D3E1}
\definecolor{bodyfill}{HTML}{E8EEF6}
\definecolor{bodyedge}{HTML}{70839D}
\definecolor{force}{HTML}{1976C9}
\definecolor{arm}{HTML}{C43E72}
\definecolor{accent}{HTML}{007F73}

\newcommand{\TitleFont}{\fontsize{40}{46}\selectfont\bfseries}
\newcommand{\PanelFont}{\fontsize{29}{35}\selectfont\bfseries}
\newcommand{\FormulaFont}{\fontsize{42}{50}\selectfont\bfseries}
\newcommand{\BodyFont}{\fontsize{22}{28}\selectfont}
\newcommand{\LabelFont}{\fontsize{23}{28}\selectfont\bfseries}

\tikzset{
  card box/.style={fill=card,draw=border,line width=1.6pt,rounded corners=18pt},
  force arrow/.style={draw=force,line width=3.4pt,-{Stealth[length=13pt,width=10pt]}},
  action line/.style={draw=force!68,line width=2pt,dash pattern=on 10pt off 7pt},
  dimension line/.style={draw=arm,line width=2.2pt},
  extension line/.style={draw=arm!78,line width=1.6pt},
  guide line/.style={draw=muted!65,line width=1.8pt,dash pattern=on 7pt off 6pt}
}

\begin{document}
\begin{tikzpicture}[x=1pt,y=1pt,font=\sffamily]
  \path[use as bounding box] (0,0) rectangle (1100,650);
  \fill[canvas] (0,0) rectangle (1100,650);
  \draw[border,line width=1.5pt,rounded corners=20pt]
    (12,12) rectangle (1088,638);

  \node[font=\TitleFont,text=ink] at (550,598)
    {Плечо и момент силы};

  \draw[card box] (42,65) rectangle (715,550);
  \draw[card box] (745,65) rectangle (1058,550);

  \node[font=\PanelFont,text=force] at (378,510)
    {Плечо силы};

  % Rigid body viewed along the rotation axis.
  \path[fill=bodyfill,draw=bodyedge,line width=2.4pt]
    (138,180) .. controls (122,302) and (170,445) .. (330,458)
    .. controls (454,468) and (594,435) .. (620,329)
    .. controls (642,238) and (547,151) .. (412,142)
    .. controls (286,154) and (155,142) .. cycle;

  % Rotation axis through O, perpendicular to the drawing plane.
  \fill[ink] (300,220) circle (5.5pt);
  \draw[ink,line width=2pt] (300,220) circle (13pt);
  \node[font=\LabelFont,text=ink,anchor=north west] at (318,204) {$O$};
  \node[font=\BodyFont,text=muted,align=left,anchor=west] at (365,205)
    {ось через $O$\\перпендикулярна рисунку};

  % Point of application and force line.
  \draw[action line] (86,410) -- (681,410);
  \fill[force] (535,410) circle (6pt);
  \node[font=\LabelFont,text=force,anchor=south] at (535,425) {$A$};
  \draw[force arrow] (535,410) -- (681,410);
  \node[font=\FormulaFont,text=force,anchor=south] at (635,429) {$\vec F$};
  \node[font=\BodyFont,text=force,anchor=south west] at (132,448)
    {линия действия силы};

  % Perpendicular arm OH and its engineering-style dimension.
  \coordinate (H) at (300,410);
  \draw[guide line] (300,220) -- (H);
  \fill[arm] (H) circle (5pt);
  \node[font=\LabelFont,text=arm,anchor=north east] at (288,399) {$H$};
  \draw[arm,line width=2.2pt] (300,388) -- (322,388) -- (322,410);

  % The upper extension is the force line itself. The lower extension
  % begins with a small overrun and stops short of the axis symbol.
  \draw[extension line] (86,220) -- (281,220);
  \draw[dimension line,-{Stealth[length=11pt,width=8pt]}]
    (98,286) -- (98,220);
  \draw[dimension line,-{Stealth[length=11pt,width=8pt]}]
    (98,344) -- (98,410);
  \node[font=\FormulaFont,text=arm,fill=card,rounded corners=5pt,
    inner xsep=8pt,inner ysep=3pt] at (98,315) {$l$};

  \node[font=\BodyFont,text=muted,align=center] at (378,96)
    {$OH$ — кратчайшее расстояние до линии действия силы};

  % Formula and meaning card.
  \node[font=\PanelFont,text=accent] at (901,505) {Модуль момента};
  \node[font=\FormulaFont,text=ink] at (901,416) {$\displaystyle |M|=Fl$};

  \draw[border,line width=1.5pt] (785,358) -- (1017,358);
  \node[font=\BodyFont,text=ink,align=left,anchor=west] at (785,309)
    {$F$ — модуль силы};
  \node[font=\BodyFont,text=ink,align=left,anchor=west] at (785,250)
    {$l$ — плечо силы};

  \draw[accent,line width=2pt,rounded corners=10pt]
    (780,116) rectangle (1022,205);
  \node[font=\BodyFont,text=accent,align=center] at (901,161)
    {если линия действия\\проходит через ось,\\то $l=0$ и $|M|=0$};
\end{tikzpicture}
\end{document}
